\documentclass[main.tex]{subfiles}
\begin{document}

\chapter{Expectation, Variance and Cauchy--Schwarz}\label{ch:beklenen-deger}

Knowing a whole distribution is often unnecessary or impossible; instead we work with summary numbers such as ``centre'' and ``spread''. The expectation\index{expectation@expectation} is an integral with respect to the probability measure, and it is the common language of estimation theory (unbiasedness, MSE), limit theorems (LLN, CLT\index{central limit theorem clt@central limit theorem (CLT)}) and Bayesian decisions (posterior expected loss). In this chapter we discuss the measure-theoretic definition of expectation, the basic inequalities and the notion of conditional expectation\index{conditional expectation@conditional expectation}, which will be used throughout the rest of the book.

\section{Definitions}

\begin{definition}[Expectation as a Lebesgue Integral]\label{def:beklenen}\index{expectation@expectation|textbf}
The expectation is defined in three steps:
\begin{enumerate}
  \item \emph{Simple RV:} for $X=\sum_{i=1}^k a_i\I_{A_i}$ ($A_i\in\F$ disjoint), $\E[X]=\sum_{i=1}^k a_i\P(A_i)$.
  \item \emph{Non-negative RV:} $\E[X]=\sup\{\E[Y]: Y \text{ simple},\ 0\le Y\le X\}$ (the value may be $+\infty$).
  \item \emph{General RV:} write $X=X^+-X^-$ ($X^+=\max(X,0)$, $X^-=\max(-X,0)$); if at least one of $\E[X^+]$ and $\E[X^-]$ is finite, $\E[X]=\E[X^+]-\E[X^-]$. If $\E|X|<\infty$, $X$ is called \emph{integrable} ($X\in L^1$).
\end{enumerate}
The result of these three steps is denoted by $\E[X] = \int_\Omega X \, d\P$.
\end{definition}

\begin{remark}[Computational Formulas]
By the change-of-variables formula known as the Law of the Unconscious Statistician (LOTUS\index{law of the unconscious statistician lotus@law of the unconscious statistician (LOTUS)}), for measurable $g$
\[
\E[g(X)]=\int_\R g(x)\,d\P_X(x)=\begin{cases}\sum_x g(x)p_X(x), & X \text{ discrete}\\[2pt] \int_{-\infty}^\infty g(x)f_X(x)\,dx, & X \text{ continuous}\end{cases}
\]
(provided the sum/integral converges absolutely). For $X\ge0$ the tail formula $\E[X]=\int_0^\infty \P(X>t)\,dt$ also holds.
\end{remark}

\begin{definition}[Moments, Variance, Covariance, Correlation]\index{correlation@correlation|textbf}\index{covariance@covariance|textbf}\index{variance@variance|textbf}
For $k\in\N$, $\mu_k'=\E[X^k]$ is the \emph{$k$-th raw moment} and $\mu_k=\E[(X-\E X)^k]$ the \emph{$k$-th central moment}. Furthermore,
\begin{gather*}
\Var(X) = \E[(X - \E X)^2] = \E[X^2] - (\E X)^2,\\
\Cov(X,Y)=\E[(X-\E X)(Y-\E Y)]=\E[XY]-\E X\,\E Y,
\end{gather*}
\[
\Corr(X,Y)=\rho_{XY}=\frac{\Cov(X,Y)}{\sqrt{\Var(X)\Var(Y)}} .
\]
If $\Cov(X,Y)=0$, $X$ and $Y$ are called \emph{uncorrelated}. The skewness\index{skewness@skewness} is $\gamma_1=\mu_3/\sigma^3$ and the excess kurtosis\index{kurtosis@kurtosis} is $\gamma_2=\mu_4/\sigma^4-3$.
\end{definition}

\begin{definition}[Moment Generating Function]
If $M_X(t)=\E[e^{tX}]$ is finite in a neighbourhood of $0$, it is called the \emph{moment generating function\index{moment generating function mgf@moment generating function (MGF)}} (MGF) of $X$. In that case $\E[X^k]=M_X^{(k)}(0)$, and the MGF determines the distribution uniquely. (Not every distribution has an MGF; this gap will be filled by the characteristic function\index{characteristic function@characteristic function} in Chapter~\ref{ch:clt}.)
\end{definition}

\begin{definition}[Conditional Expectation --- the Elementary Case]\label{def:kosullu-beklenen}\index{expectation@expectation|textbf}
Let $(X,Y)$ be a random vector with joint density (or pmf) $f_{X,Y}$, and let $f_Y(y)>0$. The conditional density of $X$ given $Y=y$ is $f_{X\mid Y}(x\mid y)=f_{X,Y}(x,y)/f_Y(y)$, and with
\[
m(y)=\E[X\mid Y=y]=\int x\,f_{X\mid Y}(x\mid y)\,dx
\]
the \emph{random variable} $\E[X\mid Y]=m(Y)$ is called the conditional expectation\index{conditional expectation@conditional expectation} of $X$ given $Y$. Similarly, $\Var(X\mid Y)=\E[(X-\E[X\mid Y])^2\mid Y]$.
\end{definition}

\begin{remark}[The General Definition]
In measure-theoretic terms, $\E[X\mid\mathcal{G}]$ is the (a.s.\ unique) RV that is measurable with respect to a sub-$\sigma$-algebra $\mathcal{G}\subseteq\F$ and satisfies $\E[X\I_G]=\E[\E[X\mid\mathcal G]\I_G]$ for every $G\in\mathcal{G}$; its existence follows from the Radon--Nikodym theorem. $\E[X\mid Y]=\E[X\mid\sigma(Y)]$. In $L^2$, $\E[X\mid\mathcal G]$ is the \emph{orthogonal projection} of $X$ onto the subspace of $\mathcal G$-measurable RVs; this geometric interpretation is the essence of the Rao--Blackwell\index{rao blackwell theorem@Rao--Blackwell theorem} theorem (Chapter~\ref{ch:yeterlilik}).
\end{remark}

\section{Main Theorems and Proofs}

\begin{proposition}[Properties of Expectation]\label{prop:beklenen-ozellik}\index{expectation@expectation}
Let $X,Y\in L^1$ and $a,b\in\R$.
\begin{enumerate}
  \item (Linearity) $\E[aX+bY]=a\E X+b\E Y$.
  \item (Monotonicity) If $X\le Y$ a.s., then $\E X\le\E Y$.
  \item If $X\indep Y$, then $\E[XY]=\E X\,\E Y$; hence independent RVs are uncorrelated.
  \item $\Var(aX+b)=a^2\Var(X)$ and $\Var(X+Y)=\Var X+\Var Y+2\Cov(X,Y)$.
  \item In general, $\Var\big(\sum_{i=1}^n X_i\big)=\sum_i\Var X_i+2\sum_{i<j}\Cov(X_i,X_j)$.
\end{enumerate}
\end{proposition}
\begin{proof}[Proof sketch]
(1)--(2) follow directly for simple RVs and, in general, by passing to the limit with the Monotone Convergence Theorem. (3) First for indicators, $\E[\I_A\I_B]=\P(A\cap B)=\P(A)\P(B)$, then by linearity and a limiting argument. (4)--(5) follow by expanding the definition and using linearity. For measure-theoretic details see \cite[Sections~1.4--1.6]{durrett} or \cite[Sections~15--16 and 21]{billingsley}.
\end{proof}

\begin{theorem}[Cauchy--Schwarz Inequality]\label{thm:cauchy-schwarz}\index{cauchy schwarz inequality@Cauchy--Schwarz inequality}
For any two random variables with $\E[X^2] < \infty$ and $\E[Y^2] < \infty$,
\[
(\E[XY])^2 \le \E[X^2]\,\E[Y^2].
\]
Equality holds if and only if $aX+bY=0$ a.s.\ for some $(a,b)\neq(0,0)$.
\end{theorem}
\begin{proof}
First, $\E[XY]$ is finite because $|XY|\le (X^2+Y^2)/2$. If $\E[X^2]=0$, then $X=0$ a.s.\ and both sides vanish. Let $\E[X^2]>0$. For every $t \in \R$,
\[
g(t) = \E[(tX + Y)^2] = t^2\E[X^2] + 2t\E[XY] + \E[Y^2] \ge 0 .
\]
A quadratic polynomial with positive leading coefficient that is never negative cannot have a positive discriminant:
$\Delta = (2\E[XY])^2 - 4\E[X^2]\E[Y^2] \le 0 \implies (\E[XY])^2 \le \E[X^2]\E[Y^2]$.
In the case of equality $\Delta=0$, so there is a $t_0$ with $g(t_0)=\E[(t_0X+Y)^2]=0$, i.e.\ $t_0X+Y=0$ a.s.
\end{proof}

\begin{corollary}[Bounds on the Correlation]\label{cor:korelasyon}\index{correlation@correlation}
$|\Corr(X,Y)|\le1$, with equality if and only if $Y=a+bX$ a.s.\ ($b\ne0$).
\end{corollary}
\begin{proof}
Apply Theorem~\ref{thm:cauchy-schwarz} to the centred variables $X-\E X$ and $Y-\E Y$.
\end{proof}

\begin{theorem}[Jensen's Inequality]\label{thm:jensen}\index{jensen s inequality@Jensen's inequality}
If $\varphi:\R\to\R$ is convex, $X\in L^1$ and $\varphi(X)\in L^1$, then
\[
\varphi(\E X)\le\E[\varphi(X)].
\]
If $\varphi$ is strictly convex, equality holds only when $X$ is a.s.\ constant.
\end{theorem}
\begin{proof}
Let $\mu=\E X$. A convex function has a supporting line at every point: $\varphi(x)\ge\varphi(\mu)+c(x-\mu)$ for all $x$ (for some $c\in\R$). Setting $x=X$ and taking expectations gives $\E\varphi(X)\ge\varphi(\mu)+c(\E X-\mu)=\varphi(\mu)$.
\end{proof}

\begin{theorem}[Tower Rule (Law of Total Expectation)]\label{thm:kule}\index{expectation@expectation}
If $X\in L^1$, then $\E\big[\E[X\mid Y]\big]=\E[X]$. Moreover, for every bounded measurable $h$, $\E[h(Y)X\mid Y]=h(Y)\E[X\mid Y]$ (``taking out what is known'').
\end{theorem}
\begin{proof}[Proof (density case)]
\[
\E[m(Y)]=\int\Big(\int x f_{X\mid Y}(x\mid y)dx\Big)f_Y(y)dy=\iint x f_{X,Y}(x,y)\,dx\,dy=\E[X],
\]
where the order of integration was exchanged by Fubini's theorem. The second property holds because $h(Y)=h(y)$ is constant given $Y=y$.
\end{proof}

\begin{theorem}[Law of Total Variance]\label{thm:toplam-varyans}\index{variance@variance}
If $X\in L^2$, then
\[
\Var(X)=\E\big[\Var(X\mid Y)\big]+\Var\big(\E[X\mid Y]\big).
\]
\end{theorem}
\begin{proof}
Let $m(Y)=\E[X\mid Y]$. Write $X-\E X=(X-m(Y))+(m(Y)-\E X)$ and square. For the cross term the tower rule\index{tower property@tower property} gives
\[
\E\big[(X-m(Y))(m(Y)-\E X)\big]=\E\Big[(m(Y)-\E X)\,\E[X-m(Y)\mid Y]\Big]=0,
\]
because $\E[X-m(Y)\mid Y]=0$. The remaining terms are $\E[(X-m(Y))^2]=\E[\Var(X\mid Y)]$ and $\E[(m(Y)-\E X)^2]=\Var(m(Y))$ (the latter uses $\E[m(Y)]=\E X$).
\end{proof}

\begin{corollary}[Best Predictor]\label{cor:en-iyi-kestirici}
If $X\in L^2$, then among all measurable functions $g$ of $Y$, the one minimising $\E[(X-g(Y))^2]$ is $g(Y)=\E[X\mid Y]$.
\end{corollary}
\begin{proof}
The same decomposition as in the proof of the law of total variance\index{law of total variance@law of total variance} gives $\E[(X-g(Y))^2]=\E[(X-m(Y))^2]+\E[(m(Y)-g(Y))^2]$; the second term vanishes for $g=m$.
\end{proof}

\section{Worked Examples}

\begin{example}[The Binomial Distribution --- the Indicator Trick]
$X\sim\text{Bin}(n,p)$ can be written as $X=\sum_{i=1}^n Y_i$ with $Y_i\iid\text{Bernoulli}(p)$. By linearity $\E X=np$; by independence $\Var X=\sum\Var Y_i=np(1-p)$. This is much shorter than evaluating the sum $\sum_k k\binom nk p^k(1-p)^{n-k}$ directly.
\end{example}

\begin{example}[Moments of the Exponential Distribution]
For $X\sim\text{Exp}(\lambda)$, integration by parts gives $\E[X^k]=k!/\lambda^k$. Hence $\E X=1/\lambda$ and $\Var X=2/\lambda^2-1/\lambda^2=1/\lambda^2$. The MGF\index{moment generating function mgf@moment generating function (MGF)} is $M_X(t)=\lambda/(\lambda-t)$, $t<\lambda$; its derivatives give the same moments.
\end{example}

\begin{example}[A Distribution Without an Expectation: Cauchy]\label{ex:cauchy-beklenen}
For the standard Cauchy density $f(x)=\frac{1}{\pi(1+x^2)}$,
\[
\E[X^+]=\int_0^\infty\frac{x}{\pi(1+x^2)}dx=\frac{1}{2\pi}\ln(1+x^2)\Big|_0^\infty=\infty,
\]
and similarly $\E[X^-]=\infty$; that is, $\E X$ is \emph{undefined}. Although the density is symmetric about zero, one cannot say ``the mean is zero''. In Chapters~\ref{ch:esitsizlikler} and~\ref{ch:clt} this distribution will show why the assumptions of the law of large numbers and the central limit theorem\index{central limit theorem clt@central limit theorem (CLT)} are needed.
\end{example}

\begin{example}[Uncorrelated but Dependent]\label{ex:iliskisiz-bagimli}
Let $X\sim\Normal(0,1)$ and $Y=X^2$. Then $\Cov(X,Y)=\E[X^3]-\E X\,\E[X^2]=0$, yet $Y$ is a deterministic function of $X$. Correlation\index{correlation@correlation} measures only \emph{linear} dependence. (Conversely, if $(X,Y)$ is \emph{jointly} normal, uncorrelatedness is equivalent to independence: Theorem~\ref{thm:mvn}. Normal marginals alone are not enough: Example~\ref{ex:normal-degil}.)
\end{example}

\begin{example}[Applications of Jensen]
\begin{itemize}
  \item $\varphi(x)=x^2$: $(\E X)^2\le\E X^2$, i.e.\ $\Var X\ge0$.
  \item $\varphi(x)=1/x$ ($x>0$): $\E[1/X]\ge1/\E X$. For instance, computing an ``average time'' from an average speed systematically underestimates the time.
  \item $\varphi(x)=-\ln x$: $\E[\ln X]\le\ln\E X$, the probabilistic form of the arithmetic--geometric mean inequality. It will be used in Chapter~\ref{ch:mle-asimptotik} to prove the consistency of the MLE (the information inequality\index{information inequality@information inequality}).
  \item For $S=\sqrt{S^2}$: $\E[S]=\E\sqrt{S^2}<\sqrt{\E S^2}=\sigma$, so $S$ is a (downward) biased estimator of $\sigma$ (Chapter~\ref{ch:nokta-tahmin}).
\end{itemize}
\end{example}

\begin{example}[Total Variance: a Poisson--Gamma Mixture]\label{ex:poisson-gamma}
Let the annual number of claims of a policyholder be $N\mid\Lambda\sim\text{Poisson}(\Lambda)$ with risk level $\Lambda\sim\text{Gamma}(\alpha,\beta)$ (shape $\alpha$, rate $\beta$). By the tower rule\index{tower property@tower property} and the law of total variance\index{law of total variance@law of total variance}\index{variance@variance},
\[
\E N=\E[\Lambda]=\frac\alpha\beta,\qquad
\Var N=\E[\Lambda]+\Var(\Lambda)=\frac\alpha\beta+\frac{\alpha}{\beta^2}.
\]
Since $\Var N>\E N$, the mixture produces \emph{overdispersion\index{overdispersion@overdispersion}} relative to the Poisson. The marginal distribution of $N$ is negative binomial. This structure is the probabilistic counterpart of the Gamma--Poisson conjugacy\index{conjugate prior@conjugate prior} in Chapter~\ref{ch:bayes}.
\end{example}

\section{Computational Applications}

\subsection{Verifying the Cauchy--Schwarz Inequality}

\pycode{code/ch03/c03_1_cauchy_schwarz_check.py}
\begin{lstlisting}[style=output]
(E[XY])^2: 9.5445 <= E[X^2]E[Y^2]: 79.7080
Linear relation: (E[XcYc])^2 = 63.7241 , E[Xc^2]E[Yc^2] = 63.7241
Sample correlation (for rho=0.85): 0.8489
\end{lstlisting}

Theoretical values: $\E[XY]=\Cov+\E X\E Y=5.1-2=3.1$, hence $(\E XY)^2=9.61$; $\E X^2=4+4=8$, $\E Y^2=9+1=10$, product $80$. The inequality is rather loose because Cauchy--Schwarz\index{cauchy schwarz inequality@Cauchy--Schwarz inequality} was applied to \emph{raw} moments; in centred form the ratio is $\rho^2=0.7225$.

\subsection{Uncorrelatedness versus Independence; the Law of Total Variance}

\pycode{code/ch03/c03_2_uncorrelated_vs_independent_total_variance.py}
\begin{lstlisting}[style=output]
Corr(X, X^2) = -0.0006   (uncorrelated)
P(Y>1 | |X|>1) = 1.0  vs  P(Y>1) = 0.3171  (dependent)
E N   : MC=6.005  theory=6.000
Var N : MC=18.034  theory=18.000  (would be 6.0 for a Poisson)
\end{lstlisting}

\subsection{Jensen's Inequality: $\E[1/X]$ and $\E[S]$}

\pycode{code/ch03/c03_3_jensen_inverse_and_s.py}
\begin{lstlisting}[style=output]
E[1/X] = 0.2501 >= 1/E[X] = 0.2000  (theory E[1/X]=1/4)
E[S] MC = 1.8787 < sigma = 2.0  | theory c4*sigma = 1.8800
\end{lstlisting}

Here $c_4(n)=\sqrt{2/(n-1)}\,\Gamma(n/2)/\Gamma((n-1)/2)$ is the correction factor giving $\E[S]=c_4\sigma$ for a normal sample ($c_4\approx0.940$ for $n=5$).

\section{Comments and Notes}

\begin{remark}[Numerical Stability When Computing a Variance]
The formula $\Var X=\E X^2-(\E X)^2$ is correct in theory but numerically dangerous: if the mean is large relative to the standard deviation, one subtracts two large numbers and floating-point errors are magnified (catastrophic cancellation). In practice a two-pass algorithm or Welford's one-pass algorithm should be used; \texttt{np.var} uses the two-pass method.
\end{remark}

\begin{remark}[Hierarchy of Moment Existence]
For $1\le p<q$, $\E|X|^q<\infty\implies\E|X|^p<\infty$ (Lyapunov's inequality $(\E|X|^p)^{1/p}\le(\E|X|^q)^{1/q}$, a consequence of Jensen\index{jensen s inequality@Jensen's inequality}). For example, the $t_\nu$ distribution has moments only of order less than $\nu$; $t_2$ has a mean but infinite variance\index{variance@variance}.
\end{remark}

\begin{remark}[Links Ahead]
The Cauchy--Schwarz\index{cauchy schwarz inequality@Cauchy--Schwarz inequality} inequality is the key step for the Cramér--Rao\index{cramer rao lower bound crlb@Cramér--Rao lower bound (CRLB)} lower bound in Chapter~\ref{ch:cramer-rao}; the law of total variance\index{law of total variance@law of total variance} for the Rao--Blackwell\index{rao blackwell theorem@Rao--Blackwell theorem} theorem in Chapter~\ref{ch:yeterlilik}; and Jensen's inequality for the proof of the consistency of the MLE in Chapter~\ref{ch:mle-asimptotik}.
\end{remark}

% --- CHAPTER SUMMARY ---
\section{Chapter Summary}
\begin{summary}
\begin{itemize}
  \item The expectation is a Lebesgue integral (Definition~\ref{def:beklenen}). It is linear and monotone, and multiplicative for independent variables (Proposition~\ref{prop:beklenen-ozellik}).
  \item The Cauchy--Schwarz inequality (Theorem~\ref{thm:cauchy-schwarz}) yields $|\Corr(X,Y)|\le1$, with equality only for a linear relation (Corollary~\ref{cor:korelasyon}).
  \item Jensen's inequality: if $\varphi$ is convex, $\varphi(\E X)\le\E\varphi(X)$ (Theorem~\ref{thm:jensen}). The information inequality and the Rao--Blackwell theorem are applications of it.
  \item The tower rule $\E[\E(X\mid Y)]=\E X$ (Theorem~\ref{thm:kule}) and the law of total variance $\Var X=\E\Var(X\mid Y)+\Var\E(X\mid Y)$ (Theorem~\ref{thm:toplam-varyans}) simplify moment calculations in hierarchical models.
  \item $\E(X\mid Y)$ is the predictor that minimises the mean squared error among functions of $Y$ (Corollary~\ref{cor:en-iyi-kestirici}). This is the basic idea behind regression and Rao--Blackwell.
  \item The moment generating function produces moments through its derivatives. When it is not defined in a neighbourhood of zero, the characteristic function of Chapter~\ref{ch:clt} is used instead.
\end{itemize}
\end{summary}
% --- END OF CHAPTER SUMMARY ---

\paragraph{Further reading.} \cite[Chapter~2]{casella}, \cite[Chapter~3]{wasserman}; for expectation\index{expectation@expectation} as a Lebesgue integral see \cite[Sections~1.5--1.6]{durrett}.

\section{Exercises}

\begin{exercise}\label{exr:03-1}
Show using Fubini's theorem that $\E X=\int_0^\infty\P(X>t)\,dt$ for $X\ge0$. Use this to find $\E X=1/p$ for $X\sim\text{Geometric}(p)$ (on $\{1,2,\dots\}$).
\hfill\hyperref[coz:03-1]{\footnotesize$\rightarrow$\,Solution}
\end{exercise}

\begin{exercise}\label{exr:03-2}
Let $X\sim\text{Uniform}(0,1)$ and $Y\mid X=x\sim\text{Uniform}(0,x)$. Compute $\E Y$ and $\Var Y$ using the tower rule and the law of total variance, and verify the result by Monte Carlo.
\hfill\hyperref[coz:03-2]{\footnotesize$\rightarrow$\,Solution}
\end{exercise}

\begin{exercise}\label{exr:03-3}
Show that $\Var(X)=\min_{c\in\R}\E[(X-c)^2]$ and $\E|X-m|=\min_{c}\E|X-c|$ ($m$ a median). How do these two results relate to the Bayes estimators of Chapter~\ref{ch:bayes}?
\hfill\hyperref[coz:03-3]{\footnotesize$\rightarrow$\,Solution}
\end{exercise}

\begin{exercise}\label{exr:03-4}
Let $Z_1,Z_2\iid\Normal(0,1)$, $X=Z_1$ and $Y=\rho Z_1+\sqrt{1-\rho^2}Z_2$. Show that $\Corr(X,Y)=\rho$, $\E[Y\mid X]=\rho X$ and $\Var(Y\mid X)=1-\rho^2$.
\hfill\hyperref[coz:03-4]{\footnotesize$\rightarrow$\,Solution}
\end{exercise}

\chapterlinks{03}

\end{document}
